% GENERATED FILE. Edit canonical lessons, then run npm run generate:books.
% canonical-source-hash: fnv1a64:250eb6f057899917
% canonical-lessons: JA-C507-kibou, JA-C507-nozomu, JA-C507-shibou-suru, JA-C507-risou, JA-C507-hoshigaru

\chapter{Hopes and Wishes}
\label{ch:ja-hopes-and-wishes}
\hlchaptermodality{507}

\begin{chapteropening}
\textbf{By the end of this chapter:} \emph{I can say what I would do if a wish came true.}

Complete the last lesson of chapter 507: I can say what I would do if a wish came true.
\end{chapteropening}

\section[\texorpdfstring{\ja{きぼう} (kibō)}{きぼう (kibō)}]{\ja{きぼう} (kibō) — a hope; a wish}

\label{lesson:JA-C507-kibou}



\begin{quote}
Before the new one: say the Japanese for a defeat, then the Japanese for to defeat.
\end{quote}

\subsection*{You'll want to know: \ja{きぼう}}
\textbf{\ja{きぼう}} — \emph{kibō} — \textquotedblleft{}a hope; a wish\textquotedblright{}.

\begin{quote}
\textbf{\ja{もし} \ja{きぼうが} \ja{あったら、おしえて} \ja{ください。}} — \emph{moshi kibō ga attara, oshiete kudasai} — \textquotedblleft{}If you have any wishes, please tell me.\textquotedblright{}
\end{quote}

\begin{cousinweb}
\textbf{\ja{きぼう}} (\emph{kibō}) is Sino-Japanese (kango): two Chinese characters read with their on'yomi, readings borrowed from Chinese. They are \emph{ki}, \textquotedblleft{}hope\textquotedblright{} (Middle Chinese \emph{xj+j}) and \emph{bō}, \textquotedblleft{}ambition\textquotedblright{} (\emph{mjangH}) (KANJIDIC2; Middle Chinese in Baxter's transcription).
\end{cousinweb}

\begin{grammarlens}[title={What you've built}]
One more word for saying what happens if.
\end{grammarlens}

\subsection*{Your turn}
\begin{itemize}
\raggedright
  \item \emph{Say it:} \emph{kibō}
  \item \emph{Say it:} \emph{kibō}, once more
  \item \emph{Say it:} \emph{moshi kibō ga attara, oshiete kudasai} — If you have any wishes, please tell me.
  \item \emph{Recall:} say the Japanese for a defeat, then the Japanese for to defeat, then say \emph{kibō} again
\end{itemize}

\begin{tcolorbox}[breakable,colback=teal!4,colframe=teal!35!black,title={Before you move on}]
What does \ja{きぼう} mean? (\textquotedblleft{}A hope; a wish\textquotedblright{}.) Say it once more.
\end{tcolorbox}
\section[\texorpdfstring{\ja{のぞむ} (nozomu)}{のぞむ (nozomu)}]{\ja{のぞむ} (nozomu) — to hope for; to desire}

\label{lesson:JA-C507-nozomu}



\begin{quote}
Before the new one: say the Japanese for to defeat, then the Japanese for a hope.
\end{quote}

\subsection*{You'll want to know: \ja{のぞむ}}
\textbf{\ja{のぞむ}} — \emph{nozomu} — \textquotedblleft{}to hope for; to desire\textquotedblright{}.

\begin{quote}
\textbf{\ja{もし} \ja{かれが} \ja{のぞんだら、この} \ja{しごとを} \ja{まかせます。}} — \emph{moshi kare ga nozondara, kono shigoto o makasemasu} — \textquotedblleft{}If he wants it, I will leave this job to him.\textquotedblright{}
\end{quote}

Polite \textbf{\ja{のぞみます}}, \emph{nozomimasu}; a godan verb. Past (polite): \textbf{\ja{かれが} \ja{ちいさい} \ja{うちを} \ja{のぞみました。}}, \emph{kare ga chiisai uchi o nozomimashita}.

\begin{cousinweb}
\textbf{\ja{のぞむ}} (\emph{nozomu}) is a native Japanese word (wago), inherited rather than borrowed from Chinese; its kanji was attached later as a kun'yomi, and its older origin is not settled (JMdict, KANJIDIC2).
\end{cousinweb}

\begin{grammarlens}[title={What you've built}]
One more word for saying what happens if.
\end{grammarlens}

\subsection*{Your turn}
\begin{itemize}
\raggedright
  \item \emph{Say it:} \emph{nozomu}
  \item \emph{Say it:} \emph{nozomu}, once more
  \item \emph{Say it:} \emph{moshi kare ga nozondara, kono shigoto o makasemasu} — If he wants it, I will leave this job to him.
  \item \emph{Recall:} say the Japanese for to defeat, then the Japanese for a hope, then say \emph{nozomu} again
\end{itemize}

\begin{tcolorbox}[breakable,colback=teal!4,colframe=teal!35!black,title={Before you move on}]
What does \ja{のぞむ} mean? (\textquotedblleft{}To hope for; to desire\textquotedblright{}.) Say it once more.
\end{tcolorbox}
\section[\texorpdfstring{\ja{しぼうする} (shibō suru)}{しぼうする (shibō suru)}]{\ja{しぼうする} (shibō suru) — to aspire to; to apply for (a school)}

\label{lesson:JA-C507-shibou-suru}



\begin{quote}
Before the new one: say the Japanese for a hope, then the Japanese for to hope for.
\end{quote}

\subsection*{You'll want to know: \ja{しぼうする}}
\textbf{\ja{しぼうする}} — \emph{shibō suru} — \textquotedblleft{}to aspire to; to apply for (a school)\textquotedblright{}.

\begin{quote}
\textbf{\ja{もし} \ja{この} \ja{がっこうを} \ja{しぼうしたら、しけんが} \ja{あるでしょう。}} — \emph{moshi kono gakkō o shibōshitara, shiken ga aru deshō} — \textquotedblleft{}If you apply for this school, there will probably be an exam.\textquotedblright{}
\end{quote}

Polite \textbf{\ja{しぼうします}}, \emph{shibōshimasu}; a \ja{する} verb. Past (polite): \textbf{\ja{むすめが} \ja{いしゃを} \ja{しぼうしました。}}, \emph{musume ga isha o shibōshimashita}.

\begin{cousinweb}
\textbf{\ja{しぼうする}} (\emph{shibō suru}) is the Sino-Japanese (kango) noun \emph{shibō} with the native verb \emph{suru}, \textquotedblleft{}to do\textquotedblright{}: two Chinese characters read with their on'yomi, readings borrowed from Chinese. They are \emph{shi}, \textquotedblleft{}intention\textquotedblright{} and \emph{bō}, \textquotedblleft{}ambition\textquotedblright{} (Middle Chinese \emph{mjangH}) (KANJIDIC2; Middle Chinese in Baxter's transcription).
\end{cousinweb}

\begin{grammarlens}[title={What you've built}]
One more word for saying what happens if.
\end{grammarlens}

\subsection*{Your turn}
\begin{itemize}
\raggedright
  \item \emph{Say it:} \emph{shibō suru}
  \item \emph{Say it:} \emph{shibō suru}, once more
  \item \emph{Say it:} \emph{moshi kono gakkō o shibōshitara, shiken ga aru deshō} — If you apply for this school, there will probably be an exam.
  \item \emph{Recall:} say the Japanese for a hope, then the Japanese for to hope for, then say \emph{shibō suru} again
\end{itemize}

\begin{tcolorbox}[breakable,colback=teal!4,colframe=teal!35!black,title={Before you move on}]
What does \ja{しぼうする} mean? (\textquotedblleft{}To aspire to; to apply for (a school)\textquotedblright{}.) Say it once more.
\end{tcolorbox}
\section[\texorpdfstring{\ja{りそう} (risō)}{りそう (risō)}]{\ja{りそう} (risō) — an ideal}

\label{lesson:JA-C507-risou}



\begin{quote}
Before the new one: say the Japanese for to hope for, then the Japanese for to aspire to.
\end{quote}

\subsection*{You'll want to know: \ja{りそう}}
\textbf{\ja{りそう}} — \emph{risō} — \textquotedblleft{}an ideal\textquotedblright{}.

\begin{quote}
\textbf{\ja{もし} \ja{りそうの} \ja{うちが} \ja{みつかったら、かいます。}} — \emph{moshi risō no uchi ga mitsukattara, kaimasu} — \textquotedblleft{}If I find my ideal house, I will buy it.\textquotedblright{}
\end{quote}

\begin{cousinweb}
\textbf{\ja{りそう}} (\emph{risō}) is Sino-Japanese (kango): two Chinese characters read with their on'yomi, readings borrowed from Chinese. They are \emph{ri}, \textquotedblleft{}logic\textquotedblright{} (Middle Chinese \emph{liX}) and \emph{sō}, \textquotedblleft{}concept\textquotedblright{} (\emph{sjangX}) (KANJIDIC2; Middle Chinese in Baxter's transcription).
\end{cousinweb}

\begin{grammarlens}[title={What you've built}]
One more word for saying what happens if.
\end{grammarlens}

\subsection*{Your turn}
\begin{itemize}
\raggedright
  \item \emph{Say it:} \emph{risō}
  \item \emph{Say it:} \emph{risō}, once more
  \item \emph{Say it:} \emph{moshi risō no uchi ga mitsukattara, kaimasu} — If I find my ideal house, I will buy it.
  \item \emph{Recall:} say the Japanese for to hope for, then the Japanese for to aspire to, then say \emph{risō} again
\end{itemize}

\begin{tcolorbox}[breakable,colback=teal!4,colframe=teal!35!black,title={Before you move on}]
What does \ja{りそう} mean? (\textquotedblleft{}An ideal\textquotedblright{}.) Say it once more.
\end{tcolorbox}
\section[\texorpdfstring{\ja{ほしがる} (hoshigaru)}{ほしがる (hoshigaru)}]{\ja{ほしがる} (hoshigaru) — to want (said of someone else)}

\label{lesson:JA-C507-hoshigaru}



\begin{quote}
Before the new one: say the Japanese for to aspire to, then the Japanese for an ideal.
\end{quote}

\subsection*{You'll want to know: \ja{ほしがる}}
\textbf{\ja{ほしがる}} — \emph{hoshigaru} — \textquotedblleft{}to want (said of someone else)\textquotedblright{}.

\begin{quote}
\textbf{\ja{もし} \ja{こどもが} \ja{ほしがったら、かって} \ja{ください。}} — \emph{moshi kodomo ga hoshigattara, katte kudasai} — \textquotedblleft{}If the child wants it, please buy it.\textquotedblright{}
\end{quote}

Polite \textbf{\ja{ほしがります}}, \emph{hoshigarimasu}; a godan verb. Past (polite): \textbf{\ja{いもうとが} \ja{にんぎょうを} \ja{ほしがりました。}}, \emph{imōto ga ningyō o hoshigarimashita}.

\begin{cousinweb}
\textbf{\ja{ほしがる}} (\emph{hoshigaru}) is a native Japanese word (wago), inherited rather than borrowed from Chinese, built on \emph{hoshii}, \textquotedblleft{}wanted\textquotedblright{}, with the ending \emph{-garu}, \textquotedblleft{}to show signs of\textquotedblright{}; its kanji was attached later as a kun'yomi (JMdict, KANJIDIC2).
\end{cousinweb}

\begin{grammarlens}[title={What you've built}]
One more word for saying what happens if.
\end{grammarlens}

\subsection*{Your turn}
\begin{itemize}
\raggedright
  \item \emph{Say it:} \emph{hoshigaru}
  \item \emph{Say it:} \emph{hoshigaru}, once more
  \item \emph{Say it:} \emph{moshi kodomo ga hoshigattara, katte kudasai} — If the child wants it, please buy it.
  \item \emph{Recall:} say the Japanese for to aspire to, then the Japanese for an ideal, then say \emph{hoshigaru} again
\end{itemize}

\begin{tcolorbox}[breakable,colback=teal!4,colframe=teal!35!black,title={Before you move on}]
What does \ja{ほしがる} mean? (\textquotedblleft{}To want (said of someone else)\textquotedblright{}.) Say it once more.
\end{tcolorbox}
